%% kapa: exam=1 round=34 year=2023 period=2 form=A subject=accounting subject_ko=회계학 seq=9 number=49
%% source: Q-Net 감정평가사 기출문제 (cst00309) artlSeq=5211789 file=raw/1008/5211789/2023년도 제34회 감정평가사 1차시험 2교시.pdf
%% conversion: kordoc 4.24.51 -> tools/gen1.py (auto); review=proofread-sonnet; flags=none
\begin{question}{49}
20×1년 말 (주)감평의 올바른 당좌예금 금액을 구하기 위한 자료는 다음과 같다. (주)감평의 입장에서 수정 전 당좌예금계정 잔액에 가산 또는 차감해야 할 금액은?
\begin{framedtext}
(1) 수정 전 잔액 \par
\hspace{2em}은행의 당좌예금잔액증명서상 금액: ￦4,000 \par
\hspace{2em}(주)감평의 당좌예금 계정원장상 금액: ￦2,100 \par
(2) 은행과 (주)감평의 당좌예금 수정 전 잔액 차이 원인 \par
○ 20×1년 말 현재 (주)감평이 발행·기록한 수표 중 은행에서 미결제된 금액: ￦1,200 \par
○ 20×1년도 은행이 기록한 수수료 미통지 금액: ￦100 \par
○ 20×1년 말 받을어음 추심으로 당좌예금 계좌에 기록되었으나, (주)감평에 미통지된 금액: ￦1,000 \par
○ 20×1년 중 거래처로부터 받아 기록하고 추심 의뢰한 수표 중 은행으로부터 부도 통지 받은 금액: ￦200 \par
\end{framedtext}
{\small\centering\begin{tabular}{@{}cll@{\hspace{3em}}cll@{}}
 & \uline{가산할 금액} & \uline{차감할 금액} & & \uline{가산할 금액} & \uline{차감할 금액} \\[2pt]
① & ￦1,000 & ￦300 & ② & ￦1,100 & ￦200 \\[2pt]
③ & ￦1,300 & ￦1,400 & ④ & ￦1,400 & ￦100 \\[2pt]
⑤ & ￦2,200 & ￦300 & & & \\
\end{tabular}\par}
\end{question}
\answer{1}
